% Source: https://artofproblemsolving.com/community/c3h3146489_2023_bmmt_team_round__berkley_mini_math_tournament_fall
% Stable identity: aops-topic-3146489-q14
Place $A=(0,0)$, $B=(5,0)$, and $P=(-3,0)$. Write $C=(u,v)$; the circle equation gives $u^2+v^2=9$, while perpendicularity gives $(u,v)\cdot(5-u,-v)=0$, hence $5u=9$. Therefore $u=\dfrac{9}{5}$ and $v=\pm\dfrac{12}{5}$. This gives $PC=\sqrt{(\dfrac{24}{5})^2+(\dfrac{12}{5})^2}=\dfrac{12\sqrt5}{5}$. Dilation about $P$ with factor $\dfrac{8}{3}$ sends the smaller circle to the larger and $C$ to $D$. Thus $CD=(\dfrac{8}{3}-1)PC=\boxed{4\sqrt5}$.
\begin{asy}
import geometry;
usepackage("amsmath");
size(280);
pair A=(0,0),B=(5,0),P=(-3,0),C=(9/5,12/5),D=P+(8/3)*(C-P);draw(circle(A,3));draw(circle(B,8));draw(P--D^^B--C^^A--C);dot(A);dot(B);label("$A$",A,SW);label("$B$",B,SE);label("$P$",P,W);label("$C$",C,NW);label("$D$",D,NE);
\end{asy}
